\section{What remains between these results and quasi-polynomial recognition}
\label{sec:global}

The new local bounds are useful because they reduce the information transported
and sometimes avoid transport entirely.  They do not show how many source
supports recognition must try.  Even if every tested candidate is cheap,
exponentially many candidates would remain an exponential recognizer.

The following elementary accounting proposition states the required global
interface precisely.  It is a conditional composition statement, not a
verified property of the present recognizer.

\begin{proposition}[Conditional hierarchy accounting]
\label{prop:accounting}
Suppose a complete recognition procedure on an $n$-crossing input visits at
most $L(n)(g(n)+1)^{L(n)}$ hierarchy stages, where $L(n)\ge1$ and $g(n)\ge0$.
Charge all preprocessing, terminal work, failed attempts, and proof assembly
to these stages.  Suppose the total charged work at each
visited stage, including candidate selection, compressed updates, intermediate
representation, and independent checking, is at most $P(n)$.  Then its total
time is at most
\[
 L(n)(g(n)+1)^{L(n)}P(n).
\]
In particular, if $L(n)=O((\log n)^2)$, $g(n)=n^{O(1)}$, and
$P(n)=2^{O((\log n)^3)}$, the complete procedure takes
$2^{O((\log n)^3)}$ time.  The same conclusion holds with polynomial $P$.
\end{proposition}

\begin{proof}
Multiply the number of visits by the cost of a visit.  Its logarithm is at
most $\log L+L\log(g+1)+\log P$, which has the claimed order under the stated
hypotheses.  Work hidden in candidate selection or in an expanding intermediate
object must be included in $P$ for the implication to apply.
\end{proof}

The shape of the visit bound occurs in Lackenby's hierarchy analysis
\cite[Proposition~9.1]{Lackenby26}; verifying the stronger parameter hypotheses
for an implemented compressed recognizer is a separate research task.  Four
specific distinctions prevent an incorrect use of \cref{prop:accounting} here.

\begin{enumerate}
\item \textbf{A supplied vector versus a useful vector.}  Support nullity and
peeling are measured after the support has been supplied.  They do not identify
a meridian support among all admissible choices.
\item \textbf{A type inventory versus marked continuations.}  The bound
$H\le2d-a$ counts unmarked coordinate types.  A hierarchy can distinguish
boundary attachments, parallelity regions, and provenance not present in that
inventory.
\item \textbf{Small output versus cheap construction.}  A short rank witness
can still require expensive discovery on an arbitrary support; a compressed
cut can have short invariants but expensive attaching maps.
\item \textbf{A local budget versus a global budget.}  Restarting a successful
local procedure at every branch does not control the sum of its work.  A
monotone potential or a complete stage-count theorem is required.
\end{enumerate}

The diagram-exterior correspondence is already checked for the canonical
input construction.  The remaining source problem is maintaining comparable
provenance through the retriangulations and compressed cuts of a future
hierarchy.  Treating the input correspondence itself as wholly missing would
misdescribe the current repository.

\section{Further research questions and proposed theorem targets}
\label{sec:agenda}

The following programme separates near-term algorithmic work from the global
theorems that would change the asymptotic recognition result.  Each item names
an observable object or a falsifiable target.  Successful experiments may
motivate a theorem but do not replace one.

\subsection{Priority A: extend the provable local savings}

\paragraph{1. Multi-seed unit elimination with a controlled residual.}
When a one-seed closure stalls, introduce a new seed and express later exposed
coordinates as integral forms in the seeds.  The unsolved original rows then
give a residual system.  Can a seed-selection rule bound the number and width
of these forms by a structural parameter that stays small under relevant
layering and cutting operations?  A plausible target is a certified reduction
to $k$ residual variables in polynomial time with controlled coefficient bits.
One must not equate the number of greedy seeds with $d$: sparse exposure
closure is weaker than linear span closure.

\paragraph{2. Characterize and minimize single-seed stopping sets.}
View the supported matching rows as a bounded-arity incidence hypergraph.
Which connected normal supports admit a complete unit-peeling order?  Can
obstructions be characterized by a small core, and can that core be computed
without large integer elimination?  Seek families with $d=1$ but growing
unavoidable core size to delimit the shortcut's applicability.  Reordering
rows alone cannot solve every stopping-set obstruction.

\paragraph{3. Optimize a measured cost, beyond the slot sum.}
The present matroid theorem optimizes $L$, but exact decoder density and
coefficient height also cost time.  Can one choose an observer with a rigorous
bicriteria bound on slot budget and decoder sparsity?  A target could compare
to the best basis under $L+\lambda\nnz(N)$, or prove an approximation guarantee
under a sparse normal-matching structure.  The simple greedy proof does not
apply automatically to this nonadditive objective.

\paragraph{4. Make general rank compilation reusable under support edits.}
Hierarchies and candidate families often differ by a few zero coordinates.
Can a certified rank/decoder update handle a deletion or insertion of $k$
columns at cost polynomial in $k$ and the affected algebraic blocks, rather
than recompiling globally?  The verifier should check a small update against
the old certified matrix while preserving exact source support binding.

\paragraph{5. Prove a practical adaptive observer policy.}
Develop a preflight estimate comparing the old sparse $q+a$ decoder, the new
$d$ observer, and the ray gate.  A useful target is an overhead guarantee:
the dispatcher takes at most a small certified budget before reverting to the
baseline, and total time is bounded by that budget plus the selected complete
query.  The negative benchmark cases in this package provide a test set for
such a policy.  Prediction accuracy alone is not an asymptotic guarantee.

\subsection{Priority B: retain geometry and markings under compression}

\paragraph{6. A marked refinement of $H\le2d-a$.}
Specify exactly which markings a hierarchy needs on a component: boundary
arcs, attachment intervals, orientation data, and maps to the original
exterior.  Does a finite invariant of controlled bit length determine the
allowed next cuts?  Seek a bound on distinct marked states in terms of $d$
and marking complexity.  A counterexample with fixed $d$ and unbounded
necessary markings would be equally valuable.  Coordinate equality alone
is insufficient for this purpose.

\paragraph{7. Source-bound compressed cutting and gluing certificates.}
Realize existing compressed-cut mathematics as an executable source-bound
operation with an independent checker.  Lackenby's Theorem~14.4 already gives
polynomial cutting and parallelity-bundle attachment data for orientable
normal surfaces in uniform-type handle structures, in handle count and
logarithmic extended weight \cite{Lackenby26}.  The new target is a checked
representation in this architecture that carries the required source maps
through those operations.  The canonical input exterior checker is an
existing base case; later cuts require their own maps and verified hypotheses.

\paragraph{8. Stable local handle types and support complexity.}
Identify the precise class of handle structures produced by an implemented
hierarchy and prove that its operations preserve a uniform local incidence
bound.  Then determine how supported nullity, peeling cores, and marking size
change under those operations.  A theorem about arbitrary triangulations
cannot silently supply a bound for an expanding hierarchy representation.

\paragraph{9. Product-region removal with a charging argument.}
Can raw-support blocks, ray certificates, and two-sided profile types identify
parallelity/product regions that can be removed while charging every removal
to a decreasing potential?  The target is a global bound on cumulative
representation size and total local work, not just termination of each
individual removal.

\subsection{Priority C: close the global selection and hierarchy bounds}

\paragraph{10. Find a useful support without exhaustive support enumeration.}
Determine whether a meridian or hierarchy surface can always be found within
a quasi-polynomially bounded family of supports, after allowed checked
geometric preprocessing.  Possible parameters include a bounded exposure
core, support nullity, and normal boundary data.  The dense Fibonacci family
shows that few occupied quadrilateral positions are not necessary for easy
certification.  A low-nullity promise alone is not a discovery algorithm.

\paragraph{11. Quantify the limits of coherent cocycle candidates.}
The maintained minimum-span and Euler-optimal coherent cocycle family is not
a complete meridian family on arbitrary solid-torus triangulations, as proved
in \code{synthesis/coherent_obstruction.tex}.  That obstruction does not
establish failure on canonical diagram exteriors.  Formulate a controlled
enlargement that repairs identified obstructions and prove its candidate count
and bit-height bounds.  Re-running
the same optimizer with a new tie-breaker does not address a structural
incompleteness theorem.

\paragraph{12. Prove a polylogarithmic hierarchy potential.}
Seek a nonnegative potential on the complete marked state whose decrease
bounds the relevant hierarchy length by $O((\log n)^2)$, or by another
polylogarithmic function sufficient for \cref{prop:accounting}.  It must charge
all reset and branching operations.  A linear bound in a complexity parameter
that may itself grow polynomially in $n$ does not yield this conclusion.

\paragraph{13. A complete bit-complexity ledger.}
For every stage, bound the number of candidates, dimensions, coordinate
heights, decoder/certificate sizes, boundary maps, and repeated local calls.
Then compose the bounds without reinitializing the accounting at each branch.
The deliverable's source hashes and separated generation/replay timings are
small-scale infrastructure for this ledger; the missing proof is global.
For negative verdicts, seek shared certificates that exclude whole families
of failed support branches after unit substitutions.  A linear dual witness
can settle an individual convex cone, but a disjunction of quadrilateral
sectors requires evidence covering every alternative.  A certificate DAG
with shared substitution prefixes is a concrete target; a single failed
candidate is not a complete rejection argument.

\subsection{Priority D: falsification and validation}

\paragraph{14. Adversarial geometric families for overhead and core growth.}
Build genuine torus-boundary triangulations with large $d$, many supported
links, persistent exposure cores, and one-sided components.  Compare exact
costs before choosing a default policy.  Abstract sparse matrices test algebra
but cannot substitute for realizable normal-matching geometries when the
claim is geometric.

\paragraph{15. Formalize the smallest trusted lemmas first.}
The triangular unit-pivot lattice theorem and integer decoder identities are
compact candidates for a proof assistant or a very small extracted checker.
Next formalize raw-incidence block separation and the primitive ray disc
criterion.  Keep a traceability map from each formal statement to the concrete
certificate fields and finite manifold hypotheses.  This would strengthen
the independent Python checks without claiming that the present package is
already formally verified.

The immediate implementation order suggested by the evidence is an adaptive
local policy, multi-seed residual compilation, and a source-bound compressed
cut interface.  The highest-value mathematical target remains a selection
or hierarchy theorem controlling the complete marked state; local arithmetic
savings alone cannot supply it.
